
\section{Limitations}
\label{sec:limits}

We state these in full because several of them bound the claims materially.

\paragraph{Absolute fidelity is low at aggressive settings.} At $\approx\!4\times$
FLOPs reduction the outputs reach SSIM $\approx\!0.53$ against the exact output:
recognisably the same image, visibly different in detail. Correction moves a
speed/fidelity front upward; it does not make aggressive caching lossless. Users
wanting near-exact output should cache less aggressively, and we did not test
correction in that regime (our correctors are trained and evaluated at anchor
intervals $4$ and $5$ only).

\paragraph{We are not a competitive accelerator.} A training-free method,
TaylorSeer \cite{liu2025taylorseer}, reports near-lossless generation at $4.99\times$
compression on FLUX.1-dev, where our corrected variant at $4.7\times$ reaches SSIM
$0.585$ against the exact output. Our operating points were chosen to make a clean
diagnostic comparison against an equal-cost correction-free baseline, not to compete,
and a reader interested in the best available speed/quality trade-off should look
there rather than here. The natural question our result raises --- whether a learned
corrector trained with a site-balanced objective could close that gap --- is one we
cannot answer without the baselines we did not run.

\paragraph{Metric scope.} All image numbers are fidelity to the same-seed exact
output. This is the right primary target for an accelerator whose purpose is to
reproduce what the model would have produced, and it is what makes same-seed paired
tests possible, but it is \emph{not} a measure of aesthetic quality or
prompt alignment: a lower SSIM means a different image, not necessarily a worse one.
We report no FID \cite{heusel2017fid} or CLIP-score on a standard benchmark such as
MS-COCO, so we cannot claim distributional quality is preserved. This is the single
largest gap in our evaluation.

\paragraph{No re-implementation of published accelerators.} Our correction-free
baseline is uniform whole-stack step caching, the mechanism underlying FORA-style
caching \cite{selvaraju2024fora} at the coarsest granularity. We do not
re-implement TeaCache \cite{liu2024teacache}, ToCa \cite{zou2024toca} or
$\Delta$-DiT \cite{chen2024deltadit}, and we make no state-of-the-art claim.

We also deliberately do \emph{not} offer oracle or $\alpha$-blend analyses as upper
bounds on a learned corrector, though we ran them: a blend's error is a contraction of
the cache error, hence collinear and direction-safe, whereas a learned corrector's
error has arbitrary direction and is amplified by Eq.~\ref{eq:multiplicative}. At
equal $\relmse$ we measured the two to differ in the \emph{sign} of recovered
quality, so the blend curve is a reference point and not a ceiling.

Our thesis is \emph{assumed} to be orthogonal to which caching schedule one starts
from; it concerns how a
learned correction on top of any of them should be trained and selected. Whether the
inversion also afflicts adaptive-schedule methods is untested, and the mechanism
makes a specific prediction there: a schedule that equalises staleness across reuse
steps flattens the per-site target-energy profile, which is the imbalance our fix
corrects, so the inversion should \emph{weaken}. Testing that on TeaCache-style
adaptive intervals is the single most informative experiment we did not run.

\paragraph{Prompt diversity.} Prompts are templated (style $\times$ subject
$\times$ condition) and the held-out split contains only $6$ subjects unseen in training.
Splits are strictly subject-disjoint, but the templating limits the breadth of the
generalisation claim; we did not evaluate on free-form or long-form prompts.

\paragraph{Timing.} Wall clock is eager mode on one A40. Absolute speedups would
change under compilation or CUDA graphs, and our own measurements shifted by
$3$--$4\%$ between sessions on unchanged baselines, which is why every speed
comparison here is internal to a single exclusive benchmark run.

\paragraph{The proxy's validity is conditional.} The per-site median tracks
deployment for correctors trained on identical data ($12/12$) and fails across
training sets ($6/26$), for the confound identified in Sec.~\ref{sec:proxy}. It
should not be used as a general-purpose quality estimator on the strength of our
evidence.

\paragraph{Mechanism is argued, not proved.} Eq.~\ref{eq:multiplicative} is a
linearisation. We test its predictions --- the granularity inversion, the damping
requirement, and the weaker damping requirement on the backbone with the flatter
$\tau$ profile --- but we do not measure $J_b$ or its spectral radius directly, so the
mechanism should be read as the best available explanation rather than a
demonstrated one.

\paragraph{Secondary-backbone caveats.} The FLUX corrector was trained with half the
tokens per site as the primary one (a storage constraint), so its absolute numbers
are not comparable to the primary backbone's; only the within-backbone loss contrast
is. Its damping optimum is interior (we swept to $\sigma{=}2$), so unlike an earlier
draft of this work we are not reporting a boundary value.

\paragraph{Our effect estimate moved four times, always toward zero as the design
improved.} We measured the loss contrast at $+0.0217$ (48 images), $+0.0115$
(192 images), $+0.0266$ (a checkpoint-selected base-data pair, which we later found
step-mismatched) and finally $+0.0077$ (step-matched, seed-replicated,
selection-free). Each estimate lay inside its successor's interval, so none was
strictly inconsistent --- but a reader quoting the first would be wrong by a factor of
three. The general lesson is that every uncontrolled degree of freedom we removed made
the effect smaller, which is the direction one should expect and a reason to distrust
the first number any pipeline produces. It is also why we claim the ordering, which
survived all four designs, rather than the magnitude, which did not.

\paragraph{We selected checkpoints with the metric we are indicting.} Every
checkpoint in this paper is the one minimising validation $\relmse$ over training,
which by our own argument is the wrong criterion. For the union-data contrast this is
harmless --- both arms select their final step ($16$k/$16$k) --- but it is \emph{not}
harmless for the base-data contrast, whose arms landed on steps $6$k and $12$k. We
found this only when re-checking the safeguard against a headline we had moved, which
is itself the lesson: a selection rule has to be re-verified against every claim that
comes to depend on it, not once. The right protocol, which we recommend and did
not follow, is to select on a deployment-aligned statistic or at a fixed step count.

\paragraph{A normalisation bug affected one reported cell.} Our feature-normalisation
statistics were computed by strided subsampling whose stride resonated with the
$4$-periodic ordering of injection sites, sampling only $\tau\in\{1,3\}$ and making
every scale $\approx\!22\%$ low. We found and fixed this mid-project (it is now
covered by unit tests), but one checkpoint used in our earlier $48$-image analysis
predates it, which is why the $192$-image tables use post-fix checkpoints only. A post-fix rerun reaches the same $\relmse$ ($0.3189$ vs.\ $0.3203$), so the bug is
benign in residual space --- but its deployed quality differs by $0.0126$ SSIM, about
six times the seed noise, and we cannot currently explain that gap (the post-fix run
was never damping-selected, which is the likeliest but untested cause). We therefore
do not claim the bug was harmless, and we rest no conclusion on that cell.

\paragraph{Two different ``gain'' baselines appear in this paper.} Gains ``vs.\ front''
interpolate the correction-free curve at matched cost; gains ``vs.\ plain caching''
compare against the same anchor schedule without a corrector. The second is
systematically larger. We use the first for headline numbers because it is the
conservative one, and label every table accordingly. Results also draw on runs with
different image-seed pairs and, for the granularity ladder, a smaller design split
($n{=}8$) than the main tables ($n{=}48$); every caption states its $n$.

\paragraph{Effect sizes are small relative to per-image variance.} The per-image
standard deviation of the paired loss difference ($0.03$--$0.05$ SSIM) exceeds the
effect itself ($0.01$--$0.03$) --- it is a shift of a distribution, not a per-image
guarantee. Depending on the condition, between $102$ and $141$ of $192$ images improve,
not all of them. Readers wanting a per-image guarantee will not find one here.

\paragraph{Scale.} Some cells of Table~\ref{tab:main} are single training runs. We
replicated the decisive contrast three times and measured the seed noise
($\sigma \approx 0.002$ at interval $5$) precisely so that single-run differences
below $\approx\!0.006$ are not interpreted; we mark such comparisons as
inconclusive rather than reporting them as effects. An earlier draft of this work
inferred a noise floor of $0.013$ from a single confounded pair and used it to doubt
real effects; replicates, not inference from one pair, are the right instrument.

\section{Conclusion}

Learned correction of cached diffusion transformers is usually posed as a regression
problem in residual space, and evaluated there. We showed that this framing is
actively misleading: at fixed granularity, data, architecture and cost, residual
error and deployed image quality rank two correctors in opposite orders, reliably and
reproducibly. The reason is an aggregation mismatch --- residual error weights
injection sites by target energy, deployment weights them equally because each feeds
the same depth-wise propagation chain --- and it is repairable. A per-site
scale-invariant objective deploys measurably better at no cost in parameters or
latency --- $+0.008$ SSIM and $-0.009$ LPIPS in our cleanest design, both significant
under subject clustering, and positive in every condition we tested --- and a per-site relative statistic restores the correct ranking as a
candidate selection criterion. We stress the ordering rather than the magnitude: the
ordering is what contradicts residual error and it is stable across samples, seeds,
metrics and data conditions, whereas our own magnitude estimate halved when we
quadrupled the evaluation set. The same lens correctly predicts that on-policy data collection, which
residual error endorses, buys nothing on images.

The broader lesson generalises past caching. Whenever a learned module is inserted
into a frozen iterative computation, the natural local reconstruction loss inherits
whatever weighting the local error happens to have, while the deployed system
propagates every insertion equally. Reporting the local metric alongside a
deployed one is not enough --- our own prior conclusions, drawn from the local metric,
were reversed twice by the deployed one.

\paragraph{Reproducibility.} Every experiment runs on a single 48\,GB GPU. The
protocol we would ask others to keep: select injection strength on validation and
never on test; replicate the decisive contrast across training seeds and report the
seed spread next to any effect; verify the deployed low-precision path reproduces the
training path before trusting a rollout (this caught a silent \texttt{fp16}
timestep-embedding failure in our own pipeline); inspect artefacts at native
resolution; and treat residual-space error strictly as a diagnostic.

\appendix
\section{Artefact measurement detail}
\label{app:artefactsdetail}

Measured over $24$ held-out cases at one image seed, in LAB space. ``Near-Nyquist'' is
the fraction of lightness-channel FFT energy at radius $>0.8\times$ Nyquist; the ratio
is cutoff-dependent, spanning $7$--$12\times$ for cutoffs $0.75$--$0.9$. ``Extreme
chroma'' is the fraction of pixels whose $(a,b)$ deviation from the exact output exceeds
$40$.

\begin{table}[h]
\centering\small
\begin{tabular}{lrrr}
\toprule
 & near-Nyquist & vs.\ exact & \% px $>40$ \\
\midrule
exact & 0.00393 & $1.00\times$ & 0.000 \\
plain cache $i{=}5$ & 0.03502 & $8.92\times$ & 5.341 \\
$+$ corrector & 0.01706 & $4.35\times$ & 3.790 \\
plain cache $i{=}2$ & 0.00909 & $2.32\times$ & 0.842 \\
\bottomrule
\end{tabular}
\end{table}

No per-image test is reported for these; $n{=}24$ at a single seed, so they are
descriptive.
\label{app:artefacts}
